% 修改入口：本文件。按用户口述与顺序手写公式复原；疑点仅用页底小字脚注记录。
\documentclass[UTF8,11pt,a4paper]{ctexart}
\usepackage[margin=2.4cm]{geometry}
\usepackage{amsmath,amssymb}
\usepackage[bottom]{footmisc}
\usepackage[hidelinks]{hyperref}
\setlength{\parskip}{0.35em}
\renewcommand{\footnotesize}{\fontsize{8}{11}\selectfont}
\newcommand{\ip}[2]{\langle #1,#2\rangle}
\newcommand{\II}{\mathrm{II}}
\begin{document}

\section{高斯映射}

我们来介绍一下什么叫高斯映照，或者叫高斯映射。以下都设 $M$ 是一个正则的、可定向的曲面；研究局部的时候，当然都是可定向的了。
考虑它的一个单位法向量场 $N$。当点在 $M$ 上变化的时候，每个点都有一个单位法向量，这个法向量肯定落在单位球面 $S^2$ 上。这样就得到了一个映射
\[
N\colon M\longrightarrow S^2.
\]
这个从 $M$ 到 $S^2$ 的映射叫作高斯映射，也就是这个法向量映射。

考虑它的微分，对每个点 $p$ 来讲，得到一个线性映射
\[
dN_p\colon T_pM\longrightarrow T_{N(p)}S^2.
\]
这里的特殊之处在哪里呢？因为 $N$ 是一个单位法向量场，可以观察到 $T_pM$ 和 $T_{N(p)}S^2$ 实际上是平行的。所以可以把它看作一个二维线性空间到它自己的映射：
\[
dN_p\colon T_pM\longrightarrow T_pM.
\]

这个微分有一些性质，其中最重要的一个性质，可能就是它是一个自伴映射。假如一个线性映射只是一个线性空间到另一个线性空间的映射，就没法说它是自伴的；只有把它看作自己到自己的线性变换，才能说它的自伴性。这也就是这里的特殊之处。

我们只需要说明它在一组基下满足自伴性。比如这组基叫作 $\omega_1,\omega_2$，那么就是要有
\[
\ip{dN_p(\omega_1)}{\omega_2}
=\ip{\omega_1}{dN_p(\omega_2)}.
\]
取一个局部参数化 $X(u,v)$，令 $\omega_1=X_u$、$\omega_2=X_v$，就可以验证这个等式。

设 $\alpha(t)=X(u(t),v(t))$ 是一条经过 $p$ 的曲线，$\alpha(0)=p$。那么
\[
\begin{aligned}
dN_p(\alpha'(0))
&=\left.\frac{d}{dt}N(u(t),v(t))\right|_{t=0}\\
&=N_u u'(0)+N_v v'(0).
\end{aligned}
\]
可以看出来，把 $X_u$ 放进 $dN_p$ 得到的就是 $N$ 对 $u$ 的导数；把 $X_v$ 放进去，得到的就是 $N$ 对 $v$ 的导数：
\[
dN_p(X_u)=N_u,\qquad dN_p(X_v)=N_v.
\]
那么只需要证明
\[
\ip{N_u}{X_v}=\ip{X_u}{N_v}.
\]

为此，肯定要对某个恒等式进行求导。把定义写出来，就是说 $N$ 和 $X_u$、$X_v$ 的内积都是零，因为它们垂直：
\[
\ip{N}{X_u}=0,\qquad \ip{N}{X_v}=0.
\]
对第一个式子关于 $v$ 求导，对第二个式子关于 $u$ 求导，得到
\[
\begin{aligned}
\ip{N_v}{X_u}+\ip{N}{X_{uv}}&=0,\\
\ip{N_u}{X_v}+\ip{N}{X_{vu}}&=0.
\end{aligned}
\]
求导以后有一个相同的项，就是 $N$ 和 $X_{uv}$ 的内积。混合偏导数的顺序可以交换，所以算一下就得到这个等式了。

\section{第二基本形式与法曲率}

考虑这样一个二次形式：
\[
Q_p(v)=\ip{dN_p(v)}{v},\qquad v\in T_pM.
\]
就是说，把原先的 $v$ 和经过 $dN_p$ 作用以后的向量做内积。给它加一个负号，就叫作第二基本形式：
\[
\II_p(v)=-Q_p(v)=-\ip{dN_p(v)}{v}.
\]

比如现在有一条曲线 $\alpha(s)$，$\alpha(0)=p$，利用它来算一下第二基本形式。把 $N$ 限制在 $\alpha$ 上，记作 $N(s)=N(\alpha(s))$，肯定就只有一个自变量 $s$ 了。现在 $N(s)$ 和 $\alpha'(s)$ 的内积为零，对它求导，得到
\[
\ip{N(s)}{\alpha''(s)}=-\ip{N'(s)}{\alpha'(s)}.
\]
而 $\alpha'(0)$ 经过 $dN_p$ 作用后，就是 $N'(0)$。所以在零这个点，第二基本形式等于法向量和曲线的二阶导数的内积：
\[
\II_p(\alpha'(0))
=-\ip{dN_p(\alpha'(0))}{\alpha'(0)}
=\ip{N(0)}{\alpha''(0)}.
\]

不妨设这条曲线是弧长参数化的。那么它的曲率就是二阶导数的模长，设主法向量为 $n(s)$，有
\[
\alpha''(s)=k(s)n(s),\qquad |\alpha''(s)|=k(s).
\]
考虑曲线的主法向量 $n$ 和曲面的法向量 $N$ 的夹角 $\theta$。因为 $n$ 和 $N$ 都是单位向量，所以
\[
\cos\theta=\ip{n}{N}.
\]
于是
\[
\begin{aligned}
\ip{N(0)}{\alpha''(0)}
&=\ip{N(0)}{k(0)n(0)}\\
&=k(0)\ip{N(0)}{n(0)}
=k(0)\cos\theta.
\end{aligned}
\]
我们给这个 $k\cos\theta$ 取一个名字，记作 $k_n$，叫作这条曲线的法曲率：
\[
k_n=k\cos\theta=\II_p(\alpha'(0)).
\]

其实可以看到一个简单的推论：经过一点 $p$ 的曲线，假如在这一点的单位切向量相同，那么法曲率也相同。因为所有这些曲线，只要切向量相同，算出来的第二基本形式就相同。所以以后只说某个方向的法曲率就可以了。

可以把法曲率解释得再直观一点。取 $p$ 点的一个切向量 $v\ne0$，考虑让 $v$ 和 $N(p)$ 张成一个平面。这个平面和曲面截出一条曲线，就可以用这条特殊的截线去算这个方向的法曲率。

\subsection{两个例子}

举一个比较简单的例子。考虑 $z=y^4$，绕 $z$ 轴旋转一周，得到一个有点像抛物面的曲面。在原点处，怎么说明 $dN_p=0$ 呢？

这时候不需要直接把 $dN_p$ 写出来，当然直接写出来也很简单。只需要说明每个方向上的法曲率都是零。那么实际上，只要证明平面曲线 $y(x)=x^4$ 在 $x=0$ 的时候曲率为零就可以了。这是二维平面内的一条曲线，求两次导数就可以完成这个工作：
\[
\gamma(x)=(x,x^4),\quad
\gamma'(x)=(1,4x^3),\quad
\gamma''(x)=(0,12x^2),\quad
\gamma''(0)=0.
\]
所以它在原点的曲率是零，每个方向的法曲率也都是零，因此 $dN_p=0$。

再比如，考虑单位球面上的第二基本形式。对每个点 $p$，考虑刚才谈到的平面截出来的法截线。很明显，这些法截线都是半径为一的圆，相当于单位球面的赤道。半径为一的圆的曲率当然是一。

再看这个圆的主法向量，很明显它和曲面向外的法向量相反，也就是向内的单位法向量。这样它的法曲率就是
\[
k_n=1\cdot(-1)=-1.
\]
所以，对每个单位切向量 $v$，都有 $\II_p(v)=-1$。

假如换一个单位球面的定向，把 $N$ 换成相反方向，很明显第二基本形式和法曲率也都会改变符号。

\section{主曲率与有关概念}

回到 $dN_p$ 是一个自伴算子这个事实。由线性代数中的标准结论，在每一点都存在一组标准正交基 $e_1,e_2$，它们是 $dN_p$ 的特征向量。比如说，记
\[
dN_p(e_1)=-k_1e_1,\qquad dN_p(e_2)=-k_2e_2.
\]
这两个 $k_1,k_2$ 就是主曲率。

这样设出来以后，实际上就可以计算任意一个方向的法曲率。比如设
\[
v=e_1\cos\theta+e_2\sin\theta,\qquad |v|=1.
\]
那么
\[
\begin{aligned}
k_n&=-\ip{dN_p(v)}{v}\\
&=-\ip{-k_1e_1\cos\theta-k_2e_2\sin\theta}
{e_1\cos\theta+e_2\sin\theta}\\
&=k_1\cos^2\theta+k_2\sin^2\theta.
\end{aligned}
\]
把这些式子展开，就得到这个经典的公式。

进一步，可以考虑 $dN_p$ 的行列式和迹。行列式就是
\[
K=\det dN_p=k_1k_2,
\]
我们把它叫作高斯曲率。平均曲率则是
\[
H=-\frac12\operatorname{tr}dN_p=\frac12(k_1+k_2).
\]\footnote{手写第2页在迹的旁边写成了 $\operatorname{tr}dN_p=-H$，口述此处未清楚保留系数。第3页及后面的口述明确把 $H$ 定义为两个主曲率的平均数，因此这里按 $\operatorname{tr}dN_p=-2H$ 复原。}

接下来是一些衍生出来的概念。

\subsection{主方向与曲率线}

已经说过，自伴映射有两个互相垂直的特征向量方向，这两个方向称为这一点处的两个主方向。

说一条曲线是曲率线，就是说，这条曲线在每一点的切向量都在某个主方向上。

有一些退化的情况。比如，假如 $dN_p$ 是恒等映射，那么当然每个方向都是主方向，这就属于退化的情况；在这样的点，对 $\alpha'$ 的方向就没有额外要求了。

\subsection{椭圆点、双曲点、抛物点和平点}

对于某一点，还可以通过 $dN_p$ 的行列式给它赋予一些名字，基本上就是椭圆点、双曲点和抛物点。

椭圆点满足 $K>0$，双曲点满足 $K<0$。行列式为零的时候，$dN_p$ 本身不一定为零。我们把 $dN_p$ 整个为零的点叫作曲面的一个平点；抛物点则相应地要求 $K=0$，但 $dN_p\ne0$。

在椭圆点，两个主曲率有相同的符号，因为高斯曲率是两个主曲率的乘积。主曲率实际上给出了所有方向上法曲率的上界和下界，从刚才含有 $\cos^2\theta$ 和 $\sin^2\theta$ 的公式就可以明显看出来。

所以，任何一个经过该点的曲线，它的主法向量都是指向该点切平面的同一侧的。

对应地，对于一个双曲点就不是同一侧了。它的两个主曲率一正一负，有的曲线的主法向量指向一侧，另一些曲线的主法向量指向另一侧。

\subsection{渐近方向、渐近线与脐点}

还有一个渐近方向的概念，就是说这个方向对应的法曲率为零，这样的方向就是渐近方向。比如，椭圆点很明显没有渐近方向。

一条线作为渐近线，和一条线作为曲率线的意思也是类似的：就是说，它在每一点的切向量都是该点的渐近方向。

另外还有脐点这个概念。假如 $k_1=k_2$，那么就说这一点是一个脐点。

\section{局部坐标中的计算}

现在来算一下 $dN_p$ 在局部坐标中的表达式。有一个参数化 $X(u,v)$，而且这个参数化和定向相容，也就是说，定向的法向量可以写成 $X_u$ 和 $X_v$ 做叉积以后归一化的结果：
\[
N=\frac{X_u\times X_v}{|X_u\times X_v|}.
\]

仍然假设 $\alpha(t)=X(u(t),v(t))$ 是一条曲线，$\alpha(0)=p$。它在 $p$ 点的切向量是
\[
\alpha'=X_u u'+X_v v'.
\]
那么
\[
dN_p(\alpha')=N_u u'+N_v v'.
\]
因为 $N_u$ 和 $N_v$ 仍然是切空间中的元素，所以可以把它们写成
\[
\begin{aligned}
N_u&=a_{11}X_u+a_{21}X_v,\\
N_v&=a_{12}X_u+a_{22}X_v.
\end{aligned}
\]
于是
\[
dN_p(\alpha')
=(a_{11}u'+a_{12}v')X_u
+(a_{21}u'+a_{22}v')X_v.
\]
假如用向量的形式写出来，在基 $X_u,X_v$ 下就是
\[
[dN_p(\alpha')]_{(X_u,X_v)}
=\begin{pmatrix}a_{11}&a_{12}\\a_{21}&a_{22}\end{pmatrix}
\begin{pmatrix}u'\\v'\end{pmatrix}.
\]
这样 $dN_p$ 的矩阵就是由这些 $a_{ij}$ 组成的矩阵。注意，这里不要求 $X_u,X_v$ 是标准正交基，所以这个矩阵其实不一定对称，也就是 $a_{12}$ 不一定等于 $a_{21}$。

进一步，考虑 $e,f,g$ 这三个量。它们的定义是让法向量分别和 $X_{uu},X_{uv},X_{vv}$ 做内积：
\[
\begin{aligned}
e&=\ip{N}{X_{uu}}=-\ip{N_u}{X_u},\\
f&=\ip{N}{X_{uv}}=-\ip{N_u}{X_v}=-\ip{N_v}{X_u},\\
g&=\ip{N}{X_{vv}}=-\ip{N_v}{X_v}.
\end{aligned}
\]
比如 $e$ 也就是 $N_u$ 和 $X_u$ 做内积以后再乘以负一，$f,g$ 也是一样。

这三个数有什么特点呢？从上面的式子代入，很容易得到一个矩阵等式。把之前讨论过的第一基本形式的三个系数记作
\[
E=\ip{X_u}{X_u},\quad F=\ip{X_u}{X_v},\quad G=\ip{X_v}{X_v}.
\]
记
\[
A=\begin{pmatrix}a_{11}&a_{21}\\a_{12}&a_{22}\end{pmatrix},
\]
那么
\[
\begin{pmatrix}e&f\\f&g\end{pmatrix}
=-A\begin{pmatrix}E&F\\F&G\end{pmatrix}.
\]\footnote{手写第3页定义的 $A$ 是前面 $dN_p$ 的坐标矩阵的转置；这里保留这个定义。该页矩阵乘积前的负号较不清楚，按 $e=-\ip{N_u}{X_u}$ 等式补出。口述中把 $A$ 与坐标矩阵混称的地方，按这两个实际写出的矩阵区分。}

现在注意，大小写的这些系数都是从曲面本身能直接算出的量，只和 $N$、$X$ 的导数有关。而 $a_{ij}$ 是我们要求的系数。把矩阵稍微变一下，就得到 $A$ 的表达式：
\[
\begin{aligned}
A
&=-\begin{pmatrix}e&f\\f&g\end{pmatrix}
\begin{pmatrix}E&F\\F&G\end{pmatrix}^{-1}\\
&=-\frac1{EG-F^2}
\begin{pmatrix}e&f\\f&g\end{pmatrix}
\begin{pmatrix}G&-F\\-F&E\end{pmatrix}.
\end{aligned}
\]\footnote{手写第3页逆矩阵一行的非对角元看起来写成了 $F$，没有看清负号，且未写出分母。这里按上一行矩阵等式复原为 $-F$，并补出 $1/(EG-F^2)$。}
当然还可以用逆矩阵的公式，把 $a_{ij}$ 直接写出来：
\[
\begin{aligned}
a_{11}&=\frac{fF-eG}{EG-F^2},&
a_{21}&=\frac{eF-fE}{EG-F^2},\\
a_{12}&=\frac{gF-fG}{EG-F^2},&
a_{22}&=\frac{fF-gE}{EG-F^2}.
\end{aligned}
\]

这样，算高斯曲率 $K$ 的时候，当然就是算这个矩阵的行列式。所以可以写成
\[
K=\det dN_p=\det A=\frac{eg-f^2}{EG-F^2}.
\]
这里要注意大小写的区别。

那么平均曲率怎么计算呢？设 $k_1,k_2$ 是那两个主曲率。注意到 $dN_p+k_iI$ 这个线性变换不可逆，所以它的行列式只能为零：
\[
\det\begin{pmatrix}
a_{11}+k&a_{12}\\a_{21}&a_{22}+k
\end{pmatrix}=0,\qquad k=k_1,k_2.
\]
展开后得到
\[
k^2+(a_{11}+a_{22})k+\det(a_{ij})=0.
\]
可以发现，$k_1,k_2$ 就是这个二次方程的根，而且常数项其实就是高斯曲率 $K$。那么可以用韦达定理，把 $k_1+k_2$ 和 $a_{11}+a_{22}$ 联系起来：
\[
a_{11}+a_{22}=-(k_1+k_2).
\]
所以平均曲率就是
\[
H=-\frac{a_{11}+a_{22}}2=\frac{k_1+k_2}2.
\]\footnote{此处口述转写成了“等于 $a_1$ 加 $a_2$，除以 $2$”，负号不清楚。手写第3页写有 $a_{11}+a_{22}=-(k_1+k_2)$，故这里保留负号。}
这也是它叫“平均”的由来，因为它是两个主曲率的平均数。

假如把 $a_{ij}$ 的表达式代进去，就得到用大小写的 $E,F,G$ 和 $e,f,g$ 表达 $H$ 的式子：
\[
H=\frac{eG-2fF+gE}{2(EG-F^2)}.
\]
进一步，假如想解，当然也能够通过 $H$ 和 $K$ 把两个主曲率解出来，因为无非就是解一个二次方程：
\[
k^2-2Hk+K=0,\qquad
k_1,k_2=H\pm\sqrt{H^2-K}.
\]\footnote{手写第3页末行写成了 $k^2-2\operatorname{tr}dN_p\,k+K=0$。这与上面的展开式及 $H=(k_1+k_2)/2$ 不一致；这里按这两处复原为 $k^2-2Hk+K=0$。}

\section{几个特定的例子}

\subsection{全脐曲面}

下面来研究几个特定的例子。第一个就是全脐曲面。我们将说明，假如一个连通曲面 $S$ 的所有点都是脐点，那么它只能是平面或者球面的一部分。\footnote{口述说“只能是平面或者球面”。仅有连通与全脐条件时，结论是包含在某个平面或球面中，不一定覆盖整个平面或球面；这里据此补出“一部分”。}

因为每个点都是脐点，所以对每个 $q\in S$，可以写成
\[
dN_q(w)=\lambda(q)w,\qquad w\in T_qS.
\]
我们首先证明 $\lambda(q)$ 在局部是常数。假如局部上是常数，很显然也可以推广到整体，因为曲面是连通的。

现在考虑参数化 $X(u,v)$。把 $w$ 写成 $w=X_u a_1+X_v a_2$，上面的方程就可以写成
\[
N_u a_1+N_v a_2=\lambda(X_u a_1+X_v a_2).
\]
由 $a_1,a_2$ 的任意性，对应的系数肯定相等，就是说
\[
N_u=\lambda X_u,\qquad N_v=\lambda X_v.
\]
为了消去一些项，对第一个式子关于 $v$ 求偏导，对第二个式子关于 $u$ 求偏导：
\[
\begin{aligned}
N_{uv}&=\lambda_vX_u+\lambda X_{uv},\\
N_{vu}&=\lambda_uX_v+\lambda X_{vu}.
\end{aligned}
\]
因为 $N_{uv}=N_{vu}$，相减以后就得到
\[
\lambda_vX_u=\lambda_uX_v.
\]
$X_u,X_v$ 线性无关，所以 $\lambda_u=\lambda_v=0$。那么 $\lambda$ 在局部是常数，在整个连通曲面上当然也就是常数了。

假如 $\lambda=0$，整个法向量 $N$ 就是常数，这当然就相当于是平面了。

下面假设 $\lambda=a\ne0$。我们的想法是直接尝试找出一个 $Y(u,v)$，证明它的导数为零，然后证明它就是这个球面的圆心。考虑
\[
Y(u,v)=X(u,v)-\frac1\lambda N(u,v).
\]
在局部参数化下对 $u,v$ 求导就可以了：
\[
Y_u=X_u-\frac1\lambda N_u=0,\qquad
Y_v=X_v-\frac1\lambda N_v=0.
\]
它对 $u,v$ 的偏导都是零，所以它是常数；再由连通性，这个圆心在整个曲面上相同。

本来要证明它是圆心，可以对距离的平方求导。但得到这个常数以后，去算 $X-Y$，会发现其实也不用求导，因为 $N$ 的模长肯定是一：
\[
|X-Y|=\left|\frac1\lambda N\right|=\frac1{|\lambda|}.
\]\footnote{手写第4页距离写成 $1/\lambda$，口述在“平方”与“长度”之间也有反复。由于 $\lambda$ 可能为负，这里按长度复原为 $1/|\lambda|$；相应的距离平方是 $1/\lambda^2$。}
这样就说明整个曲面包含在一个球面上了。

\subsection{第二基本形式决定的局部性质}

第二个要介绍的是一个由第二基本形式决定的局部性质。考虑 $p\in M$ 是一个椭圆点，我们要说明，在局部的意义下，$p$ 周围有一个邻域 $V$，使得 $V$ 中除 $p$ 外的所有点都在 $T_pM$ 的同一侧。假如是双曲点，曲面就在两侧。这里有的时候会把 $S$ 和 $M$ 混着用，它们都是曲面的意思。

怎么证明呢？把曲面上的点到切平面的有向距离表示出来就可以了。还是选择 $X(u,v)$ 作为参数化，令 $p=X(0,0)$。我们需要 $p$ 点的法向量 $N(p)$，有向距离可以由这样的内积得到：
\[
d(u,v)=\ip{X(u,v)-X(0,0)}{N(p)}.
\]
一般都假设参数化 $X$ 是光滑的，所以可以写二元泰勒公式：
\[
\begin{aligned}
X(u,v)={}&X(0,0)+X_u u+X_v v\\
&+\frac12\bigl(X_{uu}u^2+2X_{uv}uv+X_{vv}v^2\bigr)+R(u,v),
\end{aligned}
\]
其中各阶导数取在 $(0,0)$，余项满足
\[
R(u,v)=o(u^2+v^2).
\]
把它代进去算一算，切向的一次项与 $N(p)$ 的内积为零，得到
\[
\begin{aligned}
d(u,v)&=\frac12(eu^2+2fuv+gv^2)+R_1(u,v)\\
&=\frac12\II_p(\omega)+R_1(u,v),\qquad
\omega=X_u u+X_v v,
\end{aligned}
\]
其中 $R_1=\ip{R}{N(p)}$ 仍然是一个小量。

第二基本形式是二次的，也就是说，$\omega$ 乘以一个系数，它的值就乘以这个系数的平方。又因为 $X_u,X_v$ 是一组基，$R_1=o(|\omega|^2)$。所以对 $\omega\ne0$，有
\[
\frac{d(u,v)}{|\omega|^2}
=\frac12\II_p\left(\frac{\omega}{|\omega|}\right)+o(1).
\]\footnote{手写第4页这一行看起来是 $d|\omega|^2$，口述的“齐次化”也较模糊。按前面的二次展开，这里应当除以 $|\omega|^2$，把第二基本形式的自变量归一化。}

在椭圆点，不妨选择法向，使第二基本形式是正的。归一化以后，只需要看单位切向量上第二基本形式的取值范围。它有一个严格大于零的下界，比如记作 $m>0$。

当 $(u,v)\to(0,0)$ 时，$\omega$ 也趋于零。这个时候，$o(1)$ 最终会比第二基本形式的一半的最小下界还小，所以
\[
\frac{d(u,v)}{|\omega|^2}>0.
\]
因此，这些点总是在切平面的同一侧。

刚才说的是椭圆点。双曲点也是类似的，因为第二基本形式在单位切向量上既取正值，也取负值。分别沿这两个方向趋于原点，余项足够小以后，有向距离也分别为正和为负，所以曲面在切平面的两侧都有点。

这个证明主要用了微积分中的泰勒公式，肯定要用到可微性之类的条件。\footnote{文字转写为“非积分的定理”，对应的具体用词无法确认；结合前面的二元泰勒展开，暂按“微积分中的泰勒公式”复原。}

\subsection{杜邦标线}

第三个就是杜邦标线（Dupin 标线）。仍然设 $p$ 是曲面 $M$ 中的一个点。在 $p$ 点的杜邦标线，指的是 $T_pM$ 中的子集
\[
\{\omega\in T_pM:|\II_p(\omega)|=1\}.
\]
也就是说，第二基本形式作用在 $\omega$ 上的值是 $\pm1$。

为了化简这个东西，取 $e_1,e_2$ 为两个主方向组成的标准正交基。这里不能随便取一组标准正交基，要取对应 $k_1,k_2$ 的那一组。

考虑 $\omega=\rho v$，其中 $|v|=1$，$\rho$ 就是半径。这是一个极坐标表示，可以写成
\[
v=e_1\cos\theta+e_2\sin\theta.
\]
用之前得到的公式，有
\[
\begin{aligned}
\pm1=\II_p(\omega)
&=\rho^2\II_p(v)\\
&=k_1\rho^2\cos^2\theta+k_2\rho^2\sin^2\theta\\
&=k_1x^2+k_2y^2,
\end{aligned}
\]
这里 $x=\rho\cos\theta$、$y=\rho\sin\theta$，是 $\omega$ 在 $e_1,e_2$ 下的坐标。

这实际上说明，标线的形状是二次曲线，可以对它进行分类。对于椭圆点，$k_1,k_2$ 同号，由高中的几何知识就知道它代表一个椭圆。这里，假如两个主曲率为正，$\pm1$ 就只能取 $+1$。\footnote{口述说“这里的正负一就只能取一了”，没有明确说已选取正的定向。若 $k_1,k_2<0$，应当取 $-1$；正文按两个主曲率为正的情形复原。}

那么，这个标线有什么几何上的解释呢？以下考虑 $M$ 是一个图像曲面：以 $x,y$ 为参数，$z$ 是它的高度，$z=h(x,y)$。

通过平移、旋转，可以假设所考虑的点是
\[
p=(0,0,h(0,0))=(0,0,0),\qquad h_x(0,0)=h_y(0,0)=0.
\]
也就是说，它在零处的切平面恰好是水平的。

来计算一下这种情况下的相关系数。令
\[
X(x,y)=(x,y,h(x,y)),
\]
那么
\[
X_x=(1,0,h_x),\qquad X_y=(0,1,h_y),
\]
\[
N=\frac{(-h_x,-h_y,1)}{\sqrt{1+h_x^2+h_y^2}}.
\]
再算出二阶导数：
\[
X_{xx}=(0,0,h_{xx}),\quad
X_{xy}=(0,0,h_{xy}),\quad
X_{yy}=(0,0,h_{yy}).
\]
和 $N$ 做内积，就可以算第二基本形式中的 $e,f,g$：
\[
e=\frac{h_{xx}}{\sqrt{1+h_x^2+h_y^2}},\quad
f=\frac{h_{xy}}{\sqrt{1+h_x^2+h_y^2}},\quad
g=\frac{h_{yy}}{\sqrt{1+h_x^2+h_y^2}}.
\]
当然，还可以进一步推导 $K$ 和 $H$ 的公式。

不过，对于之前说的这种情况，$h$ 对 $x,y$ 的一阶导数在零处都是零，所以在零处会有很多消去。取切向量
\[
w=aX_x+bX_y\in T_pM,
\]
用第二基本形式的标准公式，把 $e,f,g$ 代进去，就得到
\[
\II_p(w)=ea^2+2fab+gb^2
=h_{xx}(0,0)a^2+2h_{xy}(0,0)ab+h_{yy}(0,0)b^2.
\]
在微积分中知道，这实际上就是 $h$ 在原点处的 Hessian 对应的二次型。

我们想说明的是，假如 $p$ 不是平点，把 $z=0$ 这个平面向上移动一点，截出来一条曲线，也就是考虑点集
\[
C_\varepsilon=\{(x,y):h(x,y)=\varepsilon\},\qquad \varepsilon>0.
\]
假如这个点集是一条曲线，那么它近似地相似于 $M$ 在 $p$ 点的杜邦标线。\footnote{口述此处只假设“不是平点”，并取 $\varepsilon>0$。若是抛物点，二次型退化，余项可能影响截线形状；若是负定的椭圆点，向上移动也可能截不到局部曲线。这里保留后文“忽略余项”的非严格解释，不将其视为所有非平点处的严格结论；双曲点的正、负截面分别对应标线的两个符号分支。}

为此，不妨假设 $x$ 轴和 $y$ 轴的方向是主方向：$x$ 轴的单位向量是 $e_1$，$y$ 轴的单位向量是 $e_2$。计算得到
\[
k_1(p)=h_{xx}(0,0),\qquad k_2(p)=h_{yy}(0,0),\qquad h_{xy}(0,0)=0.
\]\footnote{手写第5页把 $k_1(p)$ 写成了 $\ip{N_u}{X_u}/\ip{X_u}{X_u}=h_{xx}$。按前文 $dN(e_1)=-k_1e_1$ 的约定，内积的商前应有负号，才能与后面的 $h_{xx}$ 一致。}
再把 $h$ 做泰勒展开，就可以写成
\[
h(x,y)=\frac12\bigl(k_1x^2+k_2y^2\bigr)+R(x,y),
\qquad R=o(x^2+y^2).
\]
所以曲线 $C_\varepsilon$ 近似地由方程
\[
k_1x^2+k_2y^2=2\varepsilon
\]
给出。假如加上余项，准确地说，右边应该是 $2\varepsilon-2R(x,y)$。

当然，这是不严格的解释。就是说，把 $R$ 省略掉，得到 $k_1x^2+k_2y^2=2\varepsilon$，然后做一个相似的、同时的伸缩变换，就得到杜邦标线的样子了。

\section{向量场与积分曲线}

考虑 $M$ 到 $\mathbb R^2$ 的局部坐标，以及 $M$ 上的向量场 $\omega$。取 $p\in M$，曲线 $\alpha\colon I\to U$ 满足
\[
\alpha'(t)=\omega(\alpha(t)).
\]\footnote{本节只有本次 PDF 第5页下部和第6页的手写片段，没有对应口述。开头写有 $M\to\mathbb R^2$，这里暂理解为转入局部坐标，未据此增加全局坐标的假设。下面引理及构造尚不完整，仅按可辨认的内容复原。}

一个引理：设 $\omega$ 是向量场，$\omega(p)\ne0$，取 $p$ 的邻域 $W$，存在 $f$，对 $q\in W$ 有 $df_q\ne0$。\footnote{手写引理在“存在 $f$”之后的具体映射范围、$df_q\ne0$ 是否还应含作用的向量，均无法确定；暂照可辨部分录入，不补成完整定理。}

在 $\mathbb R^2$ 中，取 $p=(0,0)$，$\omega(p)=(1,0)$。考虑
\[
\alpha\colon U\times I\longrightarrow W,\qquad
\alpha(q,0)=q.
\]
$\alpha(q,t)$ 表示从 $q$ 出发、经过时间 $t$ 的流。再限制到 $x=0$：
\[
\widetilde\alpha
=\alpha\big|_{(U\times I)\cap\{(x,y,t):x=0\}}.
\]

\section{曲面的内蕴几何：等距与共形}

下面要处理一些曲面的内蕴几何学的基本知识，比如绝妙定理、协变导数之类的东西。

\subsection{等距对应}

首先解释一下什么叫等距对应。这里仍然考虑嵌入在 $\mathbb R^3$ 里的曲面。比如有两个曲面 $S$ 和 $\overline S$，首先取它们之间的一个微分同胚
\[
\varphi\colon S\longrightarrow\overline S.
\]
当然，局部的也可以，只需要把考虑的对象限制在局部。

我们说 $\varphi$ 是等距对应，如果对于任意 $p\in S$ 以及 $\omega_1,\omega_2\in T_pS$，都有
\[
\ip{\omega_1}{\omega_2}_p
=\ip{d\varphi_p(\omega_1)}{d\varphi_p(\omega_2)}_{\varphi(p)}.
\]
意思就是说，这个微分同胚保持第一基本形式的所有信息不变。

在以后的叙述中，假如想强调等距对应是局部的，就说这是一个局部等距对应。强调“局部”的时候，通常是在说，这个等距对应不一定能延伸到整体。

圆柱面和平面是局部等距的，但是它们甚至不同胚，很显然没法讨论整体的等距对应。\footnote{这一句的文字转写只辨出“路面。平面”，手写中没有对应名称；暂按上下文复原为“圆柱面和平面”。这里说的是整个圆柱面与整个平面。}

注意，下面的判定是很明显的。对于 $S$ 和 $\overline S$，假如存在 $\mathbb R^2$ 中的一个开集 $U$，以及两个参数表示
\[
X\colon U\longrightarrow S,\qquad
\overline X\colon U\longrightarrow\overline S,
\]
而且它们对应的第一基本形式的系数相同，
\[
E=\overline E,\qquad F=\overline F,\qquad G=\overline G,
\]
那么 $\overline X\circ X^{-1}$ 就给出了这两个参数片之间的等距对应。

\subsection{悬链面与正螺面}

下面来看一个经典的例子：悬链面和正螺面是局部等距的。先说一下这两个曲面的定义。它们应该在特殊曲面里面出现，不过现在还没有提，也没关系，可以直接给出参数表示。

悬链面是一个绕 $z$ 轴旋转的旋转曲面，参数可以写成
\[
X(u,v)=(a\cosh v\cos u,\ a\cosh v\sin u,\ av),\qquad a>0.
\]
这个参数写出来以后，第一基本形式比较好算，直接算出来就是
\[
E=a^2\cosh^2v,\qquad F=0,\qquad G=a^2\cosh^2v.
\]

正螺面的一种参数表示是
\[
\overline X(\overline u,\overline v)
=(\overline v\cos\overline u,\ \overline v\sin\overline u,\ a\overline u),
\qquad 0<\overline u<2\pi,\quad -\infty<\overline v<\infty.
\]
这里用上横线来说明它和刚才的参数不一样。参数的范围也可以更大，比如让 $\overline u\in\mathbb R$，这不影响。

可以稍微把握一下正螺面的图像：直观地说，它就是在沿 $z$ 轴上升的同时，不断地拧这个面。因为 $\overline u$ 变化的时候，高度在变化，$\cos\overline u$ 和 $\sin\overline u$ 也在变化，相当于让这个面不断转动。

现在考虑下面这个变换：
\[
\overline u=u,\qquad \overline v=a\sinh v,
\qquad u\in(0,2\pi),\quad v\in\mathbb R.
\]
这时候正螺面的参数表示就变成
\[
\overline X(u,v)=(a\sinh v\cos u,\ a\sinh v\sin u,\ au).
\]
算一下它的第一基本形式系数，会发现
\[
\overline E=a^2\cosh^2v,\qquad
\overline F=0,\qquad
\overline G=a^2\cosh^2v,
\]
和之前是一样的。所以它们是局部等距的。

\subsection{共形映照}

之后还有一个有关共形映照的定义。说它是共形映照，并不一定就是等距映照，而是允许相差一个和 $p$ 有关的系数。具体来说，要求微分同胚 $\varphi$ 在每个 $p$ 点满足
\[
\ip{d\varphi_p(v_1)}{d\varphi_p(v_2)}
=\lambda^2(p)\ip{v_1}{v_2},\qquad v_1,v_2\in T_pS,
\]
其中 $\lambda(p)>0$。同样，也有局部共形映照这个概念。

共形映照的特点就是保持角度，但不一定保持长度。比如两个非零向量 $\alpha,\beta$ 的夹角满足
\[
\cos\theta=\frac{\ip{\alpha}{\beta}}{|\alpha|\,|\beta|}.
\]
当 $\alpha,\beta$ 都被映过去时，分子分母的系数会消掉：
\[
\frac{\ip{d\varphi_p(\alpha)}{d\varphi_p(\beta)}}
{|d\varphi_p(\alpha)|\,|d\varphi_p(\beta)|}
=\frac{\lambda^2(p)\ip{\alpha}{\beta}}
{\lambda(p)|\alpha|\,\lambda(p)|\beta|}
=\cos\theta.
\]
所以角度当然不变。

\section{高斯绝妙定理}

下面讨论一个比较重要的概念，就是高斯绝妙定理，这当然很有名。

以下考虑 $S$ 是一个可定向的正则曲面，参数化 $X$ 和这个定向相容，也就是说，把 $X_u,X_v$ 做叉积再归一化，得到法向量场 $N$。

注意到 $X_u,X_v,N$ 组成一个标架，可以用这个标架把想要的信息写出来，这是一个经典的方法。现在写成
\[
\begin{aligned}
X_{uu}&=\Gamma^1_{11}X_u+\Gamma^2_{11}X_v+L_1N,\\
X_{uv}&=\Gamma^1_{12}X_u+\Gamma^2_{12}X_v+L_2N,\\
X_{vv}&=\Gamma^1_{22}X_u+\Gamma^2_{22}X_v+L_3N,\\
N_u&=a_{11}X_u+a_{21}X_v,\\
N_v&=a_{12}X_u+a_{22}X_v.
\end{aligned}
\]
这里的 $\Gamma^k_{ij}$ 称为 Christoffel 符号。

回忆一下，$a_{ij}$ 之前已经得到了。那么来看一下，这些新的系数还能有什么发现。

首先，$L_1,L_2,L_3$ 其实很熟悉。比如把 $X_{uu}$ 和 $N$ 做内积以后，就得到 $L_1$；而这个内积就是小写的 $e$。所以
\[
L_1=e,\qquad L_2=f,\qquad L_3=g.
\]

为了计算 $\Gamma^k_{ij}$，把上面的前三个关系式和 $X_u,X_v$ 做内积。比如看第一个关系式。这里要注意，只用 $E,F,G$ 的值，不涉及导数的话，当然没有办法描述这些新的量。所以还要把 $E,F,G$ 对参数的偏导考虑进来，这仍然是和曲面内蕴的第一基本形式相关的量。

比如对 $E=\ip{X_u}{X_u}$ 关于 $u$ 求偏导，再用刚才的展开式写出来：
\[
E_u=2\ip{X_{uu}}{X_u}
=2\Gamma^1_{11}E+2\Gamma^2_{11}F.
\]
类似地，考虑 $X_{uu}$ 和 $X_v$ 的内积，也可以尝试用 $F$ 和它的偏导数凑出来：
\[
F_u=\ip{X_{uu}}{X_v}+\ip{X_u}{X_{uv}},\qquad
E_v=2\ip{X_u}{X_{uv}}.
\]
所以
\[
\ip{X_{uu}}{X_v}=F_u-\frac12E_v
=\Gamma^1_{11}F+\Gamma^2_{11}G.
\]
可以看到，这两个式子有两个未知量，系数矩阵的行列式是 $EG-F^2$，肯定不等于零：
\[
\begin{pmatrix}E&F\\F&G\end{pmatrix}
\begin{pmatrix}\Gamma^1_{11}\\\Gamma^2_{11}\end{pmatrix}
=\begin{pmatrix}\frac12E_u\\F_u-\frac12E_v\end{pmatrix}.
\]
因此，可以通过这些内蕴的信息直接求出 $\Gamma^1_{11},\Gamma^2_{11}$。其他 $\Gamma^k_{ij}$ 也是类似的。

那么下面怎么导出这个定理呢？这样做略显生硬，不过暂时也不想找别的更好的方法了，直接考虑
\[
(X_{uu})_v=(X_{uv})_u,\qquad
(X_{vv})_u=(X_{uv})_v,\qquad N_{uv}=N_{vu}.
\]
因为最后都把它们写成 $X_u,X_v,N$ 的分量，所以再把 $X_{uu},X_{uv}$ 之类的展开式代进去，最终肯定会得到这样的式子：
\[
\begin{aligned}
A_1X_u+B_1X_v+C_1N&=0,\\
A_2X_u+B_2X_v+C_2N&=0,\\
A_3X_u+B_3X_v+C_3N&=0.
\end{aligned}
\]
里面的系数都是有关 $E,F,G$ 以及它们的导数的函数；当然，这里也包括 $e,f,g$，因为展开式中也有 $L_1,L_2,L_3$。

由于 $X_u,X_v,N$ 线性无关，实际上可以得到九个量等于零，也就是 $A_i,B_i,C_i$ 都等于零。

比如来确认一下第一个关系中 $X_v$ 的系数为零。\footnote{口述这里说的是“$A_1$ 等于零”，但手写上方把 $A_1$ 放在 $X_u$ 前，而下面的具体公式用的是 $X_v$ 的分量。这里按实际展开所得的 $X_v$ 分量复原，对应上面的 $B_1=0$。}
全部代进去算，会得到
\[
\begin{aligned}
&\Gamma^1_{11}\Gamma^2_{12}
+\Gamma^2_{11}\Gamma^2_{22}+ea_{22}
+(\Gamma^2_{11})_v\\
&\qquad=\Gamma^1_{12}\Gamma^2_{11}
+(\Gamma^2_{12})^2+fa_{21}
+(\Gamma^2_{12})_u.
\end{aligned}
\]
先把这些 $\Gamma^k_{ij}$ 都放到一边，等号另一边想算的是 $fa_{21}-ea_{22}$。用之前得到的 $a_{ij}$ 的公式，有
\[
\begin{aligned}
fa_{21}-ea_{22}
&=f\frac{eF-fE}{EG-F^2}-e\frac{fF-gE}{EG-F^2}\\
&=\frac{-f^2E+egE}{EG-F^2}
=EK.
\end{aligned}
\]
这里就出现了高斯曲率 $K$。关于 $K$ 的这个公式，我们曾经推导过。所以
\[
\begin{aligned}
EK={}&(\Gamma^2_{11})_v-(\Gamma^2_{12})_u
+\Gamma^1_{11}\Gamma^2_{12}
+\Gamma^2_{11}\Gamma^2_{22}\\
&-\Gamma^1_{12}\Gamma^2_{11}-(\Gamma^2_{12})^2.
\end{aligned}
\]
右边只涉及第一基本形式的系数及其导数，实际上就证明了想要的结论：曲面的高斯曲率在局部等距对应下不变。

\section{协变导数}

下面继续研究一些内蕴几何的相关内容。为了保证概念直观，仍然考虑曲面所嵌入的外围空间。

首先要定义协变导数。考虑 $\omega$ 是 $S$ 的一个开集 $U$ 上的向量场，$p\in U$，$Y\in T_pS$ 是一个方向向量。取一条参数曲线
\[
\alpha\colon(-\varepsilon,\varepsilon)\longrightarrow U,
\qquad \alpha(0)=p,\qquad \alpha'(0)=Y.
\]\footnote{这次的文字与手写都停在选取曲线这一步，尚未给出协变导数的具体定义；暂不补写后面的内容。}

\end{document}
